% Generated draft: requires mathematical and source review.
\documentclass[11pt]{article}
\usepackage[margin=25mm]{geometry}
\usepackage{fontspec}
\setmainfont{Latin Modern Roman}
\usepackage{amsmath,amssymb,mathtools,unicode-math}
\setmathfont{Latin Modern Math}
\usepackage{enumitem,needspace,xurl,hyperref}
\setlength{\parindent}{0pt}
\setlength{\parskip}{5pt}
\setlength{\emergencystretch}{3em}
\newcommand{\N}{\mathbb{N}}
\newcommand{\Z}{\mathbb{Z}}
\newcommand{\Q}{\mathbb{Q}}
\newcommand{\R}{\mathbb{R}}
\newcommand{\C}{\mathbb{C}}
\newcommand{\F}{\mathbb{F}}
\newcommand{\A}{\mathbb{A}}
\newcommand{\PP}{\mathbb{P}}
\newcommand{\E}{\mathbb{E}}
\newcommand{\eps}{\varepsilon}
\newcommand{\dd}{\,\mathrm{d}}
\begin{document}
\section*{UnsolvedMath ID: 800209}
\textbf{Source problem code:} catalogue entry
\subsection{Definitions and mathematical statement}
Let \(M\subset\C^n\) be a smooth real hypersurface, let \(p_0\in M\), and let \(U\subset\C^n\) be a neighborhood of \(p_0\). A smooth real defining function \(r\) on an ambient neighborhood \(V\) satisfies
\[
M\cap V=\{z\in V:r(z)=0\},\qquad \dd r\ne0\quad\text{on }M\cap V.
\]
Its positive side is \(\{z\in V:r(z)>0\}\); the choice of sign of \(r\) is part of the question.

For a \(C^1\) complex-valued function \(f\) on \(M\cap U\), the tangential Cauchy--Riemann condition is
\[
Lf=0\qquad\text{for every smooth local section }L\text{ of }T^{0,1}M,
\]
where
\[
T_p^{0,1}M=\{L\in T_p^{0,1}\C^n:Lr=0\}.
\]
This defines the CR condition in the differentiable setting; the source does not specify the regularity class of its CR functions.

\textbf{Problem (B. Stensones, as attributed in the supplied source).} Determine conditions under which CR functions on \(M\cap U\) extend holomorphically to the positive side \(\{r>0\}\).

For a local interpretation, the candidate extension domain is
\[
D_+=\{z\in U\cap V:r(z)>0\}.
\]
An extension would be a holomorphic function on this domain with boundary values equal to the given CR function. However, the source leaves unspecified the regularity and sense of these boundary values, whether extension is required on all of \(D_+\) or only after shrinking about \(p_0\), and whether the unqualified notation \(\{r>0\}\) instead denotes the entire positive side in a larger defining neighborhood. These choices remain unresolved and are not additional hypotheses of the source statement.
\subsection{Short English statement}
Characterize when CR functions near a point of a smooth real hypersurface extend holomorphically to the side selected by \(r>0\). The source does not fix function regularity, the boundary-value interpretation, or the extent of the extension domain.
\subsection{Sources}
\begin{itemize}
\item Local UnsolvedMath registry, permanent ID \(800209\), source problem code \texttt{catalogue entry}, title: \emph{Several Complex Variables Problem List --- Problem 5}.
\item Buzzard, Fornaess, and Krantz, \emph{The Several Complex Variables Problem List} (1995), Problem 5. Supplied source URL: \url{https://arxiv.org/abs/math/9509202}. The record reports access on 2026-10-16 and extraction by \texttt{arxiv\_latex}; these are supplied provenance statements, not independently verified access or source checks.
\item The registry lists Buzzard, Fornaess, and Krantz in its proposer field, while the supplied problem text carries the attribution ``B. Stensones'' and contact \texttt{berit@math.lsa.umich.edu}. Both attribution records are retained without resolving their differing roles.
\end{itemize}
\end{document}
